% KAUST Academy / NLP Applications — one-page project proposal
% Saudi Dialect to MSA Rewriter with Fact Preservation
% Overleaf: use pdfLaTeX. Black-and-white academic design, matching earlier proposal.
\documentclass[11pt,a4paper]{article}
\usepackage[T1]{fontenc}
\usepackage[utf8]{inputenc}
\usepackage{newtxtext}
\usepackage{microtype}
\usepackage[a4paper,top=17mm,bottom=16mm,left=20mm,right=20mm]{geometry}
\usepackage{tabularx,array}
\usepackage{ragged2e}
\usepackage{hyperref}
\usepackage{setspace}
\hypersetup{hidelinks,pdfauthor={Saudi Dialect to MSA Project Team},pdftitle={Saudi Dialect to Modern Standard Arabic — Project Proposal}}
\pagenumbering{gobble}
\setlength{\parindent}{0pt}
\setlength{\parskip}{0pt}
\setstretch{1.055}
\renewcommand{\arraystretch}{1.20}
\newcommand{\smalllabel}[1]{{\sffamily\fontsize{8}{10}\selectfont\bfseries\MakeUppercase{#1}}}
\newcommand{\sectionline}[2]{\vspace{6.5pt}{\sffamily\fontsize{9.1}{11.5}\selectfont\bfseries #1\quad #2}\par\vspace{2.5pt}}
\begin{document}
\thispagestyle{empty}
\noindent\begin{tabularx}{\linewidth}{@{}Xr@{}}
\smalllabel{KAUST Academy / NLP Applications} & \smalllabel{Project Proposal | 24 November 2026} \\
\end{tabularx}
\vspace{4pt}\hrule height 1.1pt\vspace{12pt}

{\sffamily\bfseries\fontsize{20}{23}\selectfont Saudi Dialect to Modern Standard Arabic}\par
\vspace{3.5pt}
{\sffamily\fontsize{10.7}{13}\selectfont A Saudi dialect rewriter for a travel booking desk, with checks for missing details and factual changes}\par
\vspace{9pt}
\noindent\begin{tabularx}{\linewidth}{@{}>{\RaggedRight\arraybackslash}X >{\RaggedLeft\arraybackslash}X@{}}
{\smalllabel{Topic 6 -- Your Own Project}} & {\smalllabel{Option A -- LoRA Fine-tuning}}\\
\end{tabularx}
\vspace{8pt}\hrule height .45pt

\sectionline{TEAM}{Project members}
\noindent\begin{tabularx}{\linewidth}{@{}p{10mm}X p{26mm}@{}}
{\smalllabel{No.}} & {\smalllabel{Full name}} & {\smalllabel{Role}}\\[-1pt]
\hline
01 & \textbf{Abdulrahman Alharthi} & Team leader\\
02 & \textbf{Yasser Al-Hajlah} & Team member\\
03 & \textbf{Musaad Al-Qahtani} & Team member\\
04 & \textbf{Mohammed Al-Rashidi} & Team member\\
\end{tabularx}
\vspace{4pt}\hrule height .45pt

\sectionline{01}{Research idea and objective}
Staff at the fictional Darb Aljawza Travel booking desk convert customer chats from Saudi dialect into formal Modern Standard Arabic (MSA) requests for suppliers. A fluent rewrite can still be unsafe if it changes a name, date, time, amount, or reference code. This project will test whether small-model adaptation improves Saudi-to-MSA rewriting and whether a deterministic fact-preservation check reduces detectable factual alterations.

\sectionline{02}{Dataset and approved golden test}
Candidate training data are Saudi--MSA sentence pairs from Khalid Almeman's multidialectal parallel travel corpus (50,010 total records across dialects; \href{https://doi.org/10.5281/zenodo.17776617}{Zenodo: 10.5281/zenodo.17776617}) plus synthetic examples rich in factual details. The corpus's licence statements will be resolved before any redistribution. The proposed V3 golden test contains 100 travel-desk messages with key facts and missing-field labels. The team will review and freeze it before training, subject to instructor approval of the replacement for the archived V2.2 release.

\sectionline{03}{Proposed training approach}
We plan supervised \textbf{LoRA fine-tuning of Qwen3-1.7B} (under the 4B model-size limit), using assistant-only loss and QLoRA if GPU memory requires it. A carefully prompted untuned Qwen3 model will serve as the baseline. A training notebook and data-processing code are prepared, but real fine-tuning and model-based comparison are subsequent stages, not completed experimental results.

\sectionline{04}{System workflow and local demonstration}
The system first asks for missing dates or destinations for new bookings, or a missing reference for supported existing-booking requests. It then rewrites complete requests in MSA and checks supported facts such as numbers, dates, places and booking codes. On a mismatch, it will attempt one constrained repair; persistent mismatches will withhold the rewrite and display the original for review. The intended demonstration uses \textbf{Ollama} for local model serving and a right-to-left \textbf{Streamlit} interface. Pattern checks are not a guarantee of full semantic equivalence.

\sectionline{05}{Evaluation and expected outcome}
We will compare (i) prompted untuned Qwen3, (ii) LoRA-tuned Qwen3, and (iii) the tuned model with fact checking on the same frozen golden set. Measures will include human-rated meaning preservation and MSA quality, preserved facts, dialect remnants, wrong-to-right versus right-to-wrong changes, and released versus withheld outputs. Ten adversarial cases will test safety. The expected outcome is a reproducible local prototype and transparent comparative evidence, not guaranteed error-free rewriting.

\vfill
\vspace{6pt}\hrule height .7pt\vspace{4pt}
\noindent\begin{tabularx}{\linewidth}{@{}>{\RaggedRight\arraybackslash}X >{\RaggedLeft\arraybackslash}X@{}}
{\sffamily\fontsize{8.2}{10}\selectfont\bfseries DATA SOURCE \normalfont\href{https://doi.org/10.5281/zenodo.17776617}{doi.org/10.5281/zenodo.17776617}} & {\sffamily\fontsize{8.2}{10}\selectfont\bfseries GOLDEN TEST \normalfont 100 team-reviewed messages}\\
\end{tabularx}
\end{document}
